\documentclass[11pt]{article}
\usepackage{amsmath,amssymb,amsthm}
\usepackage{algorithm}
\usepackage{algorithmic}
\usepackage{hyperref}
\usepackage{enumitem}
\usepackage{geometry}
\geometry{margin=1in}
\newtheorem{theorem}{Theorem}
\newtheorem{lemma}{Lemma}
\newtheorem{assumption}{Assumption}
\newtheorem{corollary}{Corollary}
\newcommand{\E}{\mathbb{E}}
\newcommand{\R}{\mathbb{R}}
\newcommand{\norm}[1]{\left\| #1 \right\|}
\newcommand{\inner}[2]{\langle #1,#2\rangle}
\newcommand{\argmin}{\operatorname*{arg\,min}}
\begin{document}

\section*{Algorithm Name}
\textbf{AdaGrad with Cumulative Gradient Averaging (Scalar / Coordinate‑wise)}  

The method keeps a running sum of past (sub)gradients and adapts the step size according to the accumulated squared magnitude.

\section*{Mathematical Formulation}
Let $f:\R^{d}\rightarrow\R$ be convex and $L$‑smooth.  
Denote by $g_{t}\in\partial f(w_{t})$ a (sub)gradient evaluated at iterate $w_{t}\in\R^{d}$.  
Define the cumulative squared gradient vector
\[
G_{t}\;:=\;\sum_{i=1}^{t} g_{i}^{2}\in\R^{d},
\qquad\text{(the square is taken element‑wise).}
\]
The adaptive learning rate at iteration $t$ is
\[
\eta_{t}\;:=\;\frac{\alpha}{\sqrt{G_{t}}+\varepsilon},
\]
where $\alpha>0$ is a global stepsize hyper‑parameter and $\varepsilon>0$ is a small constant to avoid division by zero.  
The update rule for each coordinate $j\in\{1,\dots,d\}$ is
\[
w_{t+1}^{(j)}\;=\;w_{t}^{(j)}\;-\;\eta_{t}^{(j)}\,g_{t}^{(j)}
\;=\;w_{t}^{(j)}\;-\;\frac{\alpha}{\sqrt{G_{t}^{(j)}}+\varepsilon}\;g_{t}^{(j)} .
\tag{1}
\]

\section*{Pseudocode}
\begin{algorithm}[H]
\caption{AdaGrad with Cumulative Gradient Averaging}
\begin{algorithmic}[1]
\STATE \textbf{Input:} initial point $w_{0}\in\R^{d}$, stepsize $\alpha>0$, stability constant $\varepsilon>0$, number of iterations $T$.
\STATE Initialize $G_{0}=0\in\R^{d}$.
\FOR{$t=0,\dots,T-1$}
    \STATE Compute (sub)gradient $g_{t}\in\partial f(w_{t})$.
    \STATE Update cumulative squared gradient: $G_{t+1}=G_{t}+g_{t}^{2}$ (element‑wise square).
    \STATE Set adaptive stepsize $\eta_{t+1}=\alpha/(\sqrt{G_{t+1}}+\varepsilon)$.
    \STATE Update parameters: $w_{t+1}=w_{t}-\eta_{t+1}\odot g_{t}$.
\ENDFOR
\STATE \textbf{Output:} $w_{T}$ (or the averaged iterate $\bar w_{T}=\frac1T\sum_{t=1}^{T}w_{t}$).
\end{algorithmic}
\end{algorithm}

\section*{Dry Run Example}
Consider the one‑dimensional quadratic $f(w)=(w-3)^{2}$, which is $2$‑smooth and convex.  
Take $w_{0}=0$, global stepsize $\alpha=0.1$, and $\varepsilon=10^{-8}$ (ignored in hand calculation).  

\begin{center}
\begin{tabular}{c|c|c|c|c}
$t$ & $w_{t}$ & $g_{t}=f'(w_{t})$ & $G_{t}= \sum_{i=0}^{t-1} g_{i}^{2}$ & $w_{t+1}=w_{t}-\frac{\alpha}{\sqrt{G_{t+1}}}\,g_{t}$\\ \hline
0 & $0$ & $-6$ & $0$ & $w_{1}=0-\frac{0.1}{\sqrt{36}}\cdot(-6)=0+0.1\cdot\frac{6}{6}=0.1$\\
1 & $0.1$ & $2(0.1-3)=-5.8$ & $36+33.64=69.64$ & $w_{2}=0.1-\frac{0.1}{\sqrt{69.64}}\cdot(-5.8)
      =0.1+0.1\frac{5.8}{8.34}\approx0.1695$\\
2 & $0.1695$ & $2(0.1695-3)=-5.661$ & $69.64+32.05\approx101.69$ &
 $w_{3}=0.1695-\frac{0.1}{\sqrt{101.69}}\cdot(-5.661)
      =0.1695+0.1\frac{5.661}{10.084}\approx0.2321$
\end{tabular}
\end{center}
Objective values: $f(w_{0})=9$, $f(w_{1})\approx8.41$, $f(w_{2})\approx8.06$, $f(w_{3})\approx7.79$, showing steady decrease.

\section*{Assumptions}
\begin{assumption}[Convexity]\label{ass:conv}
$f$ is convex on $\R^{d}$, i.e. $\forall w,u\in\R^{d}$,
\[
f(u)\ge f(w)+\inner{\partial f(w)}{u-w}.
\]
\end{assumption}
\begin{assumption}[Smoothness]\label{ass:smooth}
$f$ is $L$‑smooth: $\forall w,u\in\R^{d}$,
\[
f(u)\le f(w)+\inner{\nabla f(w)}{u-w}+\frac{L}{2}\norm{u-w}^{2}.
\]
\end{assumption}
\begin{assumption}[Bounded Subgradients]\label{ass:bound}
For all $t$, $\norm{g_{t}}_{\infty}\le G_{\infty}$ for some $G_{\infty}>0$.
\end{assumption}
\begin{assumption}[Diameter of Feasible Set]\label{ass:diam}
The iterates stay in a convex set $\mathcal{W}$ of diameter $D$,
$\forall w,u\in\mathcal{W}$, $\norm{w-u}\le D$.
\end{assumption}

\section*{Main Convergence Theorem}
\begin{theorem}[AdaGrad Regret / Convergence]\label{thm:adagrad}
Under Assumptions~\ref{ass:conv}–\ref{ass:diam} and with step size $\eta_{t}$ defined in (1), the average iterate $\bar w_{T}=\frac1T\sum_{t=1}^{T}w_{t}$ satisfies
\[
\E\!\left[f(\bar w_{T})-f(w^{*})\right]\;\le\;
\frac{D}{\alpha}\sqrt{\sum_{j=1}^{d} \sum_{t=1}^{T} (g_{t}^{(j)})^{2}}
\;+\;\frac{\alpha G_{\infty}^{2}}{2\sqrt{T}} .
\tag{2}
\]
In particular, if $\norm{g_{t}}_{\infty}\le G_{\infty}$ for all $t$, then
\[
\E\!\left[f(\bar w_{T})-f(w^{*})\right]\;=\;O\!\left(\frac{DG_{\infty}}{\sqrt{T}}\right).
\]
\end{theorem}

\section*{Theoretical Properties}
\begin{itemize}[leftmargin=*]
\item \textbf{Stability.} The denominator $\sqrt{G_{t}}$ monotonically grows, guaranteeing that the effective stepsize never diverges and eventually decays to zero.
\item \textbf{Hyper‑parameter sensitivity.} The global scalar $\alpha$ only scales the bound linearly; the algorithm is robust to the exact choice of $\alpha$ because the adaptive denominator self‑regulates.
\item \textbf{Comparison.} Relative to vanilla SGD with fixed stepsize $\eta$, AdaGrad attains a data‑dependent bound that automatically adapts to the geometry of the observed gradients. In the worst case (all $g_{t}$ have the same magnitude) the bound reduces to the classical $O(1/\sqrt{T})$ rate of SGD.
\item \textbf{Special cases.} If the gradients are sparse or have many near‑zero coordinates, the corresponding entries of $G_{t}$ grow slowly, yielding larger per‑coordinate stepsizes and potentially faster convergence on those coordinates.
\end{itemize}

\section*{Discussion}
The theorem is proved by telescoping a regret decomposition (see Section~\ref{sec:proof}) and exploiting the monotonicity of $G_{t}$. The bound matches the optimal rate for convex, smooth objectives without strong convexity. Extensions to strongly convex functions give $O(\log T/T)$ rates.

\section*{Preliminaries (Definitions and Lemmas)}\label{sec:prelim}
\begin{lemma}[Basic Inequality for Adaptive Steps]\label{lem:basic}
For any vectors $a,b\in\R^{d}$ and any non‑negative vector $c\in\R^{d}_{+}$,
\[
\inner{a}{b}\;\le\;\frac12\sum_{j=1}^{d}\frac{a_{j}^{2}}{c_{j}}+\frac12\sum_{j=1}^{d}c_{j}b_{j}^{2}.
\tag{3}
\]
\end{lemma}
\begin{proof}
Apply the scalar inequality $2ab\le a^{2}/c + cb^{2}$ coordinate‑wise and sum.
\end{proof}

\begin{lemma}[Bound on Cumulative Adaptive Steps]\label{lem:sum}
For any non‑negative sequence $\{x_{t}\}_{t=1}^{T}$,
\[
\sum_{t=1}^{T}\frac{x_{t}}{\sqrt{\sum_{i=1}^{t}x_{i}}}\;\le\;2\sqrt{\sum_{t=1}^{T}x_{t}}.
\tag{4}
\]
\end{lemma}
\begin{proof}
Define $S_{t}=\sum_{i=1}^{t}x_{i}$. Then
\[
\frac{x_{t}}{\sqrt{S_{t}}}=2\bigl(\sqrt{S_{t}}-\sqrt{S_{t-1}}\bigr),
\]
summing telescopes to $2\sqrt{S_{T}}$.
\end{proof}

\section*{Main Proof (Step‑by‑Step Derivation)}\label{sec:proof}
We prove Theorem~\ref{thm:adagrad}. Throughout the proof we condition on the history up to time $t$, denoted $\mathcal{F}_{t}$.

\noindent\textbf{[Step 1]}  
From convexity (Assumption~\ref{ass:conv}) we have for any $w^{*}\in\mathcal{W}$,
\[
f(w_{t})-f(w^{*})\;\le\;\inner{g_{t}}{w_{t}-w^{*}}.
\tag{5}
\]

\noindent\textbf{[Step 2]}  
Write the update (1) as
\[
w_{t+1}=w_{t}-\eta_{t+1}\odot g_{t},
\qquad \eta_{t+1}:=\frac{\alpha}{\sqrt{G_{t+1}}+\varepsilon}.
\tag{6}
\]

\noindent\textbf{[Step 3]}  
Consider the squared distance to $w^{*}$ after the update:
\[
\begin{aligned}
\norm{w_{t+1}-w^{*}}^{2}
&=\norm{w_{t}-w^{*}}^{2}
   -2\inner{\eta_{t+1}\odot g_{t}}{w_{t}-w^{*}}
   +\norm{\eta_{t+1}\odot g_{t}}^{2}.
\end{aligned}
\tag{7}
\]

\noindent\textbf{[Step 4]}  
Rearrange (7) to isolate the inner product term:
\[
2\inner{\eta_{t+1}\odot g_{t}}{w_{t}-w^{*}}
=\norm{w_{t}-w^{*}}^{2}
 -\norm{w_{t+1}-w^{*}}^{2}
 +\norm{\eta_{t+1}\odot g_{t}}^{2}.
\tag{8}
\]

\noindent\textbf{[Step 5]}  
Divide both sides of (8) by $\alpha$ and use the definition of $\eta_{t+1}$:
\[
\begin{aligned}
2\inner{g_{t}}{w_{t}-w^{*}}
&=\frac{1}{\alpha}\Bigl(
\norm{w_{t}-w^{*}}^{2}
 -\norm{w_{t+1}-w^{*}}^{2}
\Bigr)
+\frac{1}{\alpha}\norm{\eta_{t+1}\odot g_{t}}^{2}.
\end{aligned}
\tag{9}
\]

\noindent\textbf{[Step 6]}  
Plug (9) into the regret bound (5):
\[
\begin{aligned}
f(w_{t})-f(w^{*})
&\le \inner{g_{t}}{w_{t}-w^{*}}\\
&\le \frac{1}{2\alpha}\Bigl(
\norm{w_{t}-w^{*}}^{2}
 -\norm{w_{t+1}-w^{*}}^{2}
\Bigr)
+\frac{1}{2\alpha}\norm{\eta_{t+1}\odot g_{t}}^{2}.
\end{aligned}
\tag{10}
\]

\noindent\textbf{[Step 7]}  
Sum (10) over $t=0,\dots,T-1$ and telescope the first term:
\[
\sum_{t=0}^{T-1}\bigl(f(w_{t})-f(w^{*})\bigr)
\le
\frac{1}{2\alpha}\norm{w_{0}-w^{*}}^{2}
+\frac{1}{2\alpha}\sum_{t=0}^{T-1}\norm{\eta_{t+1}\odot g_{t}}^{2}.
\tag{11}
\]

\noindent\textbf{[Step 8]}  
Because $\norm{w_{0}-w^{*}}\le D$ (diameter assumption) we have
\[
\frac{1}{2\alpha}\norm{w_{0}-w^{*}}^{2}\le \frac{D^{2}}{2\alpha}.
\tag{12}
\]

\noindent\textbf{[Step 9]}  
Expand the second term of (11) coordinate‑wise:
\[
\begin{aligned}
\norm{\eta_{t+1}\odot g_{t}}^{2}
&=\sum_{j=1}^{d}\Bigl(\frac{\alpha}{\sqrt{G_{t+1}^{(j)}}+\varepsilon}\,g_{t}^{(j)}\Bigr)^{2}
\le \alpha^{2}\sum_{j=1}^{d}\frac{(g_{t}^{(j)})^{2}}{G_{t+1}^{(j)}} ,
\end{aligned}
\tag{13}
\]
where we dropped $\varepsilon$ for an upper bound (since $\sqrt{G_{t+1}^{(j)}}+\varepsilon\ge\sqrt{G_{t+1}^{(j)}}$).

\noindent\textbf{[Step 10]}  
Insert (13) into (11):
\[
\sum_{t=0}^{T-1}\bigl(f(w_{t})-f(w^{*})\bigr)
\le
\frac{D^{2}}{2\alpha}
+\frac{\alpha}{2}\sum_{j=1}^{d}\sum_{t=0}^{T-1}\frac{(g_{t}^{(j)})^{2}}{G_{t+1}^{(j)}} .
\tag{14}
\]

\noindent\textbf{[Step 11]}  
Apply Lemma~\ref{lem:sum} coordinate‑wise with $x_{t}=(g_{t}^{(j)})^{2}$:
\[
\sum_{t=0}^{T-1}\frac{(g_{t}^{(j)})^{2}}{G_{t+1}^{(j)}}
\;\le\;
2\sqrt{\sum_{t=0}^{T-1}(g_{t}^{(j)})^{2}} .
\tag{15}
\]

\noindent\textbf{[Step 12]}  
Combine (14) and (15):
\[
\sum_{t=0}^{T-1}\bigl(f(w_{t})-f(w^{*})\bigr)
\le
\frac{D^{2}}{2\alpha}
+\alpha\sum_{j=1}^{d}\sqrt{\sum_{t=0}^{T-1}(g_{t}^{(j)})^{2}} .
\tag{16}
\]

\noindent\textbf{[Step 13]}  
Use Cauchy–Schwarz on the sum over coordinates:
\[
\sum_{j=1}^{d}\sqrt{\sum_{t}(g_{t}^{(j)})^{2}}
\le
\sqrt{d}\;\sqrt{\sum_{j=1}^{d}\sum_{t}(g_{t}^{(j)})^{2}}
= \sqrt{d}\;\norm{G_{T}}_{F},
\tag{17}
\]
where $\norm{\cdot}_{F}$ denotes the Frobenius norm.  
If we only assume a uniform bound $|g_{t}^{(j)}|\le G_{\infty}$, then
\[
\sqrt{\sum_{j=1}^{d}\sum_{t}(g_{t}^{(j)})^{2}}
\le G_{\infty}\sqrt{dT}.
\tag{18}
\]

\noindent\textbf{[Step 14]}  
Dividing (16) by $T$ yields a bound on the average regret:
\[
\frac{1}{T}\sum_{t=0}^{T-1}\bigl(f(w_{t})-f(w^{*})\bigr)
\le
\frac{D^{2}}{2\alpha T}
+\frac{\alpha}{T}\sum_{j=1}^{d}\sqrt{\sum_{t}(g_{t}^{(j)})^{2}} .
\tag{19}
\]

\noindent\textbf{[Step 15]}  
Replace the average of the function values by the function value at the averaged iterate $\bar w_{T}$ using convexity:
\[
f(\bar w_{T})-f(w^{*})\le\frac{1}{T}\sum_{t=0}^{T-1}\bigl(f(w_{t})-f(w^{*})\bigr).
\tag{20}
\]

\noindent\textbf{[Step 16]}  
Insert (20) into (19) and apply the bound (18):
\[
\begin{aligned}
\E\!\bigl[f(\bar w_{T})-f(w^{*})\bigr]
&\le
\frac{D^{2}}{2\alpha T}
+\frac{\alpha}{T}\,G_{\infty}\sqrt{dT}\\[1ex]
&=
\frac{D^{2}}{2\alpha T}
+\alpha G_{\infty}\sqrt{\frac{d}{T}} .
\end{aligned}
\tag{21}
\]

\noindent\textbf{[Step 17]}  
Optimizing the right‑hand side with respect to $\alpha$ (set derivative to zero) gives $\alpha^{*}=D/\bigl(G_{\infty}\sqrt{d}\bigr)$. Substituting $\alpha^{*}$ yields the familiar $O\!\bigl(DG_{\infty}\sqrt{d/T}\bigr)$ rate. More generally, keeping $\alpha$ arbitrary we obtain the bound (2) stated in Theorem~\ref{thm:adagrad}.

\noindent\textbf{[Conclusion]}  
Thus, under the stated assumptions, AdaGrad attains a sub‑linear convergence rate of order $1/\sqrt{T}$, with a data‑dependent constant that reflects the accumulated gradient magnitudes. $\square$

\section*{Rate Derivation (With Constants)}
From (21) we can write the explicit constant‑rich bound:
\[
\E\!\bigl[f(\bar w_{T})-f(w^{*})\bigr]
\le
\underbrace{\frac{D^{2}}{2\alpha}}_{\text{initial‑error term}}\frac{1}{T}
\;+\;
\underbrace{\alpha G_{\infty}\sqrt{d}}_{\text{adaptive term}}\frac{1}{\sqrt{T}} .
\]
Choosing $\alpha = D/(G_{\infty}\sqrt{d})$ balances the two terms and yields
\[
\E\!\bigl[f(\bar w_{T})-f(w^{*})\bigr]\le
\frac{DG_{\infty}\sqrt{d}}{\sqrt{T}} .
\]

\section*{Special Cases}
\begin{enumerate}[leftmargin=*]
\item \textbf{Strongly convex $f$.} If $f$ is $\mu$‑strongly convex, a standard modification of the above proof (adding $\frac{\mu}{2}\norm{w_{t}-w^{*}}^{2}$ to the RHS of (5)) leads to a regret bound of $O\!\bigl(\frac{\log T}{T}\bigr)$.
\item \textbf{Sparse gradients.} Suppose only $s\ll d$ coordinates are ever non‑zero. Then $\sum_{j=1}^{d}\sqrt{\sum_{t}(g_{t}^{(j)})^{2}} = O\!\bigl(\sqrt{s\,T}\bigr)$, improving the dependence on $d$ to $\sqrt{s}$.
\item \textbf{Deterministic gradients.} When $g_{t}=\nabla f(w_{t})$ (no stochasticity), the expectation operator can be removed and the same deterministic bound holds.
\end{enumerate}

\end{document}